\documentclass[a4paper]{article}

% Set margins
\usepackage[hmargin=2.5cm, vmargin=3cm]{geometry}

\frenchspacing

% Language packages
\usepackage[utf8]{inputenc}
\usepackage[T1]{fontenc}
\usepackage[magyar]{babel}

% AMS
\usepackage{amssymb,amsmath}

% Graphic packages
\usepackage{graphicx}

% Colors
\usepackage{color}
\usepackage[usenames,dvipsnames]{xcolor}

% Listings
\usepackage{listings}

% Question
\newenvironment{question}[1]
{\noindent\textcolor{OliveGreen}{$\circ$ \textit{#1}}

\smallskip

\color{Gray}

}{\bigskip}

% Task
\newenvironment{task}[1]
{\noindent\textcolor{RoyalBlue}{$\circ$ \textit{#1}}

\smallskip

\color{Gray}

}{\bigskip}

% Notification
\newenvironment{notification}[1]
{\noindent\textcolor{Peach}{$\circ$ \textit{#1}}

\smallskip

\color{Gray}

}{\bigskip}

% Problem
\newenvironment{problem}[1]
{\noindent\textcolor{OrangeRed}{$\circ$ \textit{#1}}

\smallskip

\color{Gray}

}{\bigskip}

% Solution
\newenvironment{solution}
{\color{RoyalBlue}}{\bigskip}

% Comment
\newenvironment{comment}
{\color{Green}}{\bigskip}

\lstset{% setup listings
        framexleftmargin=16pt,
        framextopmargin=6pt,
        framexbottommargin=6pt, 
        frame=single, rulecolor=\color{black}
}

\begin{document}

\begin{center}
   \large \textbf{Valószínűségszámítás gyakorlatok}
\end{center}

\vskip 1cm

%----------------------
\section{Kombinatorika}
%----------------------

\subsection{Eseményalgebra}

\begin{task}{Jelölje $A$ azt az eseményt, hogy 4 kockadobásból egyszer sem kapunk 6-ost. Mit jelent az $\overline{A}$?}
\end{task}

\begin{task}{Legyen $A$ az az esemény, hogy egy kockával dobva páros számot kapunk, $B$ hogy 4-nél kevesebbet dobunk, $C$ hogy legalább 3-mat dobunk. Mit jelent az alábbi esemény?}
\end{task}
\[
[A \setminus (B \cap C)] \cup [(A \setminus B) \setminus C]
\]

Legyen 2 lámpa. Jelölje $A$ azt az eseményt, hogy az első lámpa világít, $B$ pedig, hogy a második világít.
Mit jelentenek az alábbi halmazok?

\[
\overline{B}, A \cup B, A \cap B, \overline{A} \cap B, \overline{A} \cup \overline{B}, A \setminus B, \overline{A \cup B}, (A \setminus B) \cup (B \setminus A).
\]

\subsection{$\sigma$ algebra}

Melyek alkotnak $\sigma$-algebrát az alábbiak közül, és miért?

\begin{itemize}
\item $\Omega = \{1, 2\}, \mathcal{F} = \{\{1\}, \{2\}\}$
% Nincs benne az Omega és az üres halmaz!
\item $\Omega = \{1, 2, 3, 4\}, \mathcal{F} = \{\emptyset, \{1, 2\}, \{3, 4\}, \{1, 2, 3, 4\}\}$
% Igen, az.
\item $\Omega = \{1, 2, 3, 4\}, \mathcal{F} = \{\emptyset, \{1\}, \{2, 3\}, \Omega\}$
% A komplementer az Omega-ra vonatkozik!
\item $\Omega = \{1, 2, 3\}, \mathcal{F} = \{\emptyset, \{1\}, \{2\}, \{3\}, \Omega\}$
% Nincsenek benne az uniók!
\end{itemize}

\begin{comment}
\noindent Az $\mathcal{F} = \{A \subset \Omega \}$ egy $\sigma$-algebra, hogy ha
\begin{enumerate}
\item $\Omega \in \mathcal{F}$,
\item $A \in \mathcal{F}$ akkor $\overline{A} \in F$,
\item ha $A, B \in \mathcal{F}$, akkor $A \cup B \in \mathcal{F}$,
\item ha $A_1, A_2, \ldots \in \mathcal{F}$, akkor $A_1, A_2, \ldots \in \mathcal{F}$.
\end{enumerate}
Az $A \in \mathcal{F}$ egy esemény.

\end{comment}

\subsection{Binomiális tétel}

\[
\displaystyle
(a + b)^n = \sum_{k=0}^{n} \binom{n}{k} a^k b^{n - k}
\]

\begin{task}{
Fel kell rajzolni a Pascal háromszöget számértékekkel és binomiális együtthatókkal.
}
\end{task}

\noindent Tulajdonságok:
\begin{enumerate}
\item $\displaystyle \binom{n}{k} = \binom{n}{n - k}$
\item $\displaystyle \binom{n+1}{k+1} = \binom{n}{k+1} + \binom{n}{k}$
\item $\displaystyle \sum_{k=0}^{n} \binom{n}{k} = 2^n$
\end{enumerate}

\begin{itemize}
\item $(a + 2)^3 = ?$
\item Fejtsük ki az $(x^3 + y)^4$ kifejezést a binomiális tétel segítségével!
\end{itemize}

\subsection{Permutációk}

Elemek lehetséges sorrendjeit vizsgáljuk.
Amennyiben az elemek mind különbözőek \textit{ismétlés nélküli permutáció}ról beszélünk.
\begin{align*}
n! &= n \cdot (n - 1) \cdot \ldots \cdot 2 \cdot 1 \\
0! &= 1 \\
P_n &= n! \\
\end{align*}

\begin{task}{Hányféleképpen állíthatunk sorba 5 embert?}
\end{task}

\begin{task}{Hányféleképpen ültethetünk fel 7 embert egy körhintára (ha feltételezzük, hogy minden hely egyenértékű)?}
\end{task}

Amennyiben az elemek között előfordulnak azonosak, úgy \textit{ismétléses permutáció}ról beszélünk.
\[
P_{n}^{k_1, \ldots, k_r} = \dfrac{n!}{k_1! \cdots k_r!}
\]

\begin{task}{4 Arany, 8 ezüst és 3 ólom gyöngyöt szeretnénk felfűzni egy láncra. Mennyiféleképpen tehetjük ezt meg? Karika/körbe kötött lánc esetében mennyi lesz a változatok száma?}
\end{task}

\begin{task}{Egy 10 fős asztalhoz 5 nőt és 5 férfit szeretnénk leültetni, úgy hogy a nők és a férfiak felváltva következzenek. Mennyiféleképpen tehetjük ezt meg, hogy ha csak a szomszédsági viszonyokra vagyunk tekintettel?}
\end{task}

\subsection{Variációk}

$n$ darab elemből szeretnénk kiválasztani $k$ darabot úgy, hogy számít a sorrend.
Ismétlés/visszatevés nélküli eset.

\[
V_{n}^{k} = \dfrac{n!}{(n - k)!} = n \cdot (n - 1) \cdot \ldots \cdot (n - k + 1)
\]

\begin{task}{Az angol ábécéből hányféleképpen tudunk kivenni 5 betűt, úgy hogy számít azok sorrendje, és nem rakhatjuk vissza őket?}
\end{task}

Ismétléses/visszatevéses eset.

\[
V_{n}^{k, i} = n^k
\]

\begin{task}{Hányféleképpen színezhető ki 6 négyzet, ha 4 féle színünk van, és egy négyzet csak egy színű lehet?}
\end{task}

\subsection{Kombinációk}

$n$ darab elemből szeretnénk $k$ darab elemet kiválasztani.

Visszatevés/ismétlés nélküli eset.

\[
\displaystyle C_{n}^{k} = \binom{n}{k} = \dfrac{n!}{k! \cdot (n - k)!}
\]

\begin{task}{Van 4 különböző ajándékunk, amelyet 10 embernek szeretnénk szétosztani. (Mindenkinek maximum 1 jut, és 6 sajnos nem fog így kapni.) Hányféleképpen tehetjük ezt meg?}
Kombinációval és variáció osztva sorrenddel is megoldható.
\end{task}

Visszatevéses/ismétléses eset.

\[
\displaystyle C_{n}^{k, i} = \binom{n + k - 1}{k}
\]

\begin{task}{3 elemből húzunk visszatevéssel 8-szor. Mennyiféleképpen tehetjük ezt meg, hogy ha nem számít a húzások sorrendje?}
\end{task}

\subsection*{Összegzés}

\begin{tabular}{l|c|c|l}
& ismétlés nélküli & ismétléses & \\
\hline
permutáció & $n!$ & $\dfrac{n!}{k_1! \cdots k_r!}$ & csak sorrendet nézünk \\
\hline
variáció & $\dfrac{n!}{(n - k)!}$ & $\displaystyle\binom{n}{k}$ & kiválasztás ahol számít a sorrend \\
\hline
kombináció & $\displaystyle\binom{n}{k}$ & $\displaystyle\binom{n + k - 1}{k}$ & kiválasztás ahol nem számít a sorrend \\
\hline
\end{tabular}

\subsection{Vegyes feladatok}

\subsubsection{Betűk}

\begin{itemize}
\item Az $a, a, a, b, c, c$ betűknek hányféle sorrendje van?
\item Ezekből mennyi olyan van, amelynél a $b$ az első helyen szerepel?
\item Mennyi olyan van, amelynél az utolsó előtti helyen $a$ van?
% \item Hányféleképpen tudunk ebből kiválasztani 3 betűt, ha azok sorrendjét megkülönböztetjük, és nem tesszük vissza őket?
% \item Mennyi eset alakul, ha vissza is tehetjük őket?
\end{itemize}

\subsubsection{Konvex $n$-szög}

Hány átlója van egy konvex $n$-szögnek?

\subsubsection{Munkások}

10 munkást 10 gépre szeretnénk elosztani úgy, hogy az első 2 gépre csak 2 munkás jöhet szóba.
Mennyiféleképpen tehetjük ezt meg?

\subsubsection{Különböző fős csoportok}

9 embert hányféleképpen tudunk úgy csoportokra osztani, hogy a csoportok 2, 3 és 4 fősek legyenek?
% Az embereket tekintjük fixnek
% A csoportok azonosítóinak a lehetséges sorrendjeit keressük ismétléses permutációval

\subsubsection{Totó}

Hányféleképpen lehet kitölteni egy 13+1-es totószelvényt a $0, 1, X$ jelekkel?

\subsubsection{Három kocka}

Három dobókockával dobva hányféle dobáshármas adódhat?
Tekintsük azt az esetet, amikor figyelembe vesszük a sorrendet, és azt is, amikor nem!

\subsubsection{Selejtes gépek}

Egy dobozban 16 alkatrész van, amelyből 4 selejtes. Hányféleképpen vehetjük ki ezeket egyesével, hogy ha a selejteseket a végére szeretnénk hagyni?

\subsubsection{Számok}

Mennyi olyan 5 jegyű egész számot tudunk felírni, amelyek számjegyei 3-nál nagyobbak, de 6-nál nem nagyobbak?

\subsubsection{Autóverseny}

Ha 20-an indulnak egy autóversenyen, akkor hányféleképpen alakulhat az első 5 helyezett?

\subsubsection{Zsákbamacska}

Zsákokba rakunk fehér, vörös és fekete színű macskákat. Minden zsákba 4 macska kerül.
Hányféle zsák keletkezhet?
% Visszatevéses kombináció a színekkel 4 helyre.

\subsubsection{Golyók és rekeszek}

20 rekeszbe szeretnénk beletenni 10 golyót.
Egy rekeszbe több golyó is kerülhet.
A golyókat nem különböztetjük meg.
Hányféleképpen lehetséges ez?
% Visszatevés nélküli kombináció a rekeszekkel való cimkézésre nézve.

\subsubsection{Számsor}

Adottak 1-től 10-ig az egész számok.
\begin{itemize}
\item Mennyiféleképpen rendezhetjük sorba őket úgy, hogy ha szeretnénk, hogy az 5 és a 6 ilyen sorrendben kövesse egymást?
% 10! / 2
\item Mennyi az esetek száma, ha az 5 és a 6 közvetlenül, ilyen sorrendben egymás mellett kell legyen?
% 9 * 8!
\end{itemize}

\subsubsection{Úszóverseny}

Krisztina és Katinka elindul egy úszóversenyen, amelyen rajtuk kívül még 6-an neveztek.
\begin{itemize}
\item Mennyi féle lehet a verseny kimenetele, ha tudjuk, hogy ők biztosan az első 3 hely valamelyikén végeznek, és nem alakulhat ki holtverseny?
% Célszerű felrajzolni Krisztina és Katinka lehetséges helyezéseit, és melléírni, hogy mindegyiknél 6! lehetőség van.
\item Hányféleképpen alakulhat az első 3 hely?
% 6 további versenyző 6 helyre kerülhet be.
\end{itemize}

\subsubsection{Postaládák}

Van 7 levelünk, amelyeket 8 postaládába szeretnénk elhelyezni. Mennyi a lehetőségek száma, ha az alábbi esetek adottak?
\begin{itemize}
\item Minden levelet egyenértékűnek tekintünk, és egy ládába több is kerülhet?
\item Minden levél egyenértékű, és egybe legfeljebb egy kerülhet?
\item Minden levél különböző és egybe több is kerülhet?
\item Minden levél különböző és egybe csak egy kerülhet?
\end{itemize}

\subsubsection{Urna és cédulák}

Egy urnában 20 cédula van 1-től 20-ig megszámozva.
Kihúzunk 5 cédulát úgy, hogy minden húzás után a kihúzott cédulát visszatesszük.
Mennyi esetben lesz a kihúzott legkisebb szám nagyobb 6-nál?
% Ismétléses kombináció

\subsubsection{Kártyák}

Hányféleképpen oszthatunk ki 32 kártyát 4 játékosnak úgy, hogy minden játékos 8 kártyát kapjon?
% 32! / (8!)^4

\subsubsection{Monoton sorozatok}

Mennyi különböző monoton sorozat készíthető az $\{1, 2, \ldots, n\}$ halmaz elemeiből?
% növekvő és csökkenő sorozatokat is figyelembe kell venni!
% Egy adott sorrendet tekintünk = Nem számít a sorrend

\subsubsection{Büfé és italok}

Egy büfében 10 féle italt árulnak. Hányféleképpen választható ki belőle 6 darab?

\subsubsection{Táncpár választós}

Egy összejövetelen 16 hölgy és 11 férfi vesz részt.
Hányféleképpen lehet kiválasztani 4 párt belőlük? (Nők és férfiak alkothatnak jelen esetben párokat.)

\subsubsection{8 jegyű számok}

Mennyi 8 jegyű számot tudunk felírni, ha 2 darab 0, 4 darab 4-es és 2 darab 7-es számjegyünk áll rendelkezésre?

%--------------------------------
\section{Feltételes valószínűség}
%--------------------------------

\subsection{Klasszikus valószínűségi mező}

\subsubsection{Valószínűség bizonyítás}

Lássuk be az alábbiakat!
\begin{itemize}
\item $P(A \cap B) \leq P(A) + P(B)$
\item $P(A \cap B) \geq P(A) + P(B) - 1$
\end{itemize}

$0 \leq P(A \cup B) \leq 1$

$0 \leq P(A) + P(B) - P(A \cap B) \leq 1$

\subsubsection{Poincare formula}

Bizonyítsuk be, hogy ha teljesül, hogy $P(A) \geq 0.6$ és $P(B) \geq 0.9$, akkor $P(A \cap B) \geq 0.5$!

\[
P(A \cup B) = P(A) + P(B) - P(A \cap B)
\]

\subsubsection{3 kocka}

3 kockát dobunk fel egyszerre. Mennyi annak a valószínűsége, hogy a dobott számok összege 10?

\subsubsection{2 kocka}

Két kockát feldobva mennyi az alábbiak valószínűsége?
\begin{itemize}
\item Két egyenlő számot dobunk?
\item Két különböző számot dobunk?
\item A két dobott szám összege 7 és az egyik dobás 6-os?
\end{itemize}

\subsubsection{Piros-fehér golyós urna}

Egy urnában 13 golyó van, 5 piros és 8 fehér.
2 golyót húzunk belőle véletlenszerűen.
Mennyi az alábbiak valószínűsége?
\begin{itemize}
\item Mindkettő fehér lesz.
\item Mindkettő piros lesz.
\item Legalább az egyik golyó fehér.
\end{itemize}

\subsubsection{Selejtes termékek}

Egy dobozban 10 jó és 4 selejtes termék van.
Mennyi annak a valószínűsége, hogy a 4 selejtes a végére marad, hogy ha sorban vesszük ki őket?

\subsubsection{10 kocka, 5 hatos}

Mennyi annak a valószínűsége, hogy 10 kockával dobva 5 hatost dobunk?

\subsubsection{Termek}

Egy hallgató $p$ valószínűséggel tartózkodhat 5 terem valamelyikében.
Tegyük fel, hogy már 4-ben kerestük, de nem találtuk ott meg.
Mennyi a valószínűsége, hogy az 5-ödikben van?

\subsection{Geometriai valószínűség}

\subsubsection{Kör alakú céltábla}

Egy egység sugarú, kör alakú céltáblára lövünk.
A találat valószínűsége a céltáblán egyenletes eloszlású.
A táblát koncentrikus körökkel 8 részre szeretnénk osztani, úgy hogy minden részbe egyenlő valószínűséggel legyen találat.
Hogyan tehetjük ezt meg?

\subsubsection{Villamos}

Úticélunkat két villamossal ($a$ és $b$) tudjuk elérni.
Az $a$ 5 percenként közlekedik, a $b$ pedig 12 percenként.
Az első villamosok 5 órakor indulnak.
7 és 7:30 között véletlenszerűen érkezünk a megállóba.
Mennyi a valószínűsége, hogy a megállóban nem kell 2 percnél többet várakozni?

\subsubsection{Pontok távolságai}

Vegyünk fel két pontot a $[0, 1]$ intervallumon véletlenszerűen.
Mennyi annak a valószínűsége, hogy ezek távolsága kisebb, mint a 0-hoz közelebbi pont 0-tól való távolsága?

\subsubsection{Találkozó}

Ketten, $a$ és $b$ találkozni szeretnének éjfél és 1 óra között.
Mindketten az adott időintervallumon belül, egymástól függetlenül érkeznek.
$a$ maximum 10 percet vár a másikra, míg $b$ maximum csak 5 percet.

Mennyi a valószínűsége, hogy így fognak tudni találkozni?

\subsubsection{Valós számok összege}

Véletlenszerűen felírunk két, 1-nél kisebb pozitív számot.
Mekkora az alábbiak valószínűsége?
\begin{itemize}
\item Az összegük 1-nél kisebb.
\item A szorzatuk kisebb, mint $\dfrac{2}{9}$-ed.
\item Összegük kisebb 1-nél és szorzatuk kisebb $\dfrac{2}{9}$-ednél.
\end{itemize}

\subsection{Mintavételezés valószínűségei}

\subsubsection{Visszatevéses mintavétel}

\[
\displaystyle
P(\xi = k) = \dfrac{\binom{n}{k} \cdot s^k \cdot (N - s)^{n-k}}{N^n}
\]

\subsubsection{Visszatevés nélküli mintavétel}

\[
\displaystyle
P(\xi = k) = \dfrac{\binom{s}{k} \cdot \binom{N - s}{n - k}}{\binom{N}{n}}
\]

\subsubsection{Feladatok}

\begin{itemize}
\item Egy dobozban 18 alkatrész van. Ezek közül 4 selejtes. Mintát veszünk belőle visszatevéssel.
Mennyi a valószínűsége, hogy közben pontosan 2 selejtet vettünk ki?
\item Egy urnában 24 golyó van, amelyből 9 piros színű, a többi fehér.
Visszatevés nélkül húzunk belőle 5 golyót.
Mennyi a valószínűsége, hogy 3 fehéret és 2 pirosat húztunk ki?
\end{itemize}

\subsection{Feltételes valószínűség}

Teljes valószínűség tétele

Bayes tétel

\[
P(A|B) \geq P(A \cap B)
\]

Mikor lehetnek egyenlőek egymással?

\subsubsection{Mennyi?}

Ismertek a $P(A|B) = 0.7, P(A|\overline{B}) = 0.3, P(B|A) = 0.6$ valószínűségek.
Mivel lesz egyenlő $P(A)$?

\subsubsection{Rendezgetés}

Ismertek az alábbi valószínűségek.
\[
P(A) = \dfrac{1}{4}, \quad
P(A|B) = \dfrac{1}{4}, \quad
P(B|A) = \dfrac{1}{2}.
\]

Számítsuk ki a $P(A \cup B)$ és $P(\overline{A}|\overline{B})$ valószínűségeket!

\subsubsection{Kifejtős}

Bizonyítsuk be, hogy
\[
P(A \cap B|C) = P(A|B \cap C)P(B|C),
\]
hogy ha $P(B \cap C) \neq 0$.

\subsection{Üzem}

Egy üzembe 8 és 12 óra között 0.8 valószínűséggel érkezik meg egy szállítmány.
11-ig még nem érkezett meg.
Mennyi a valószínűsége, hogy 11 és 12 között még beérkezik, hogy ha feltételezzük, hogy az érkezés valószínűsége egyenletes eloszlású a 8-12 óra közötti időintervallumon?

\subsection{Kérdéses válasz}

Felteszünk valakinek egy kérdést, amelyre 3 alternatív válaszlehetőség van.
Az illető $0.5$ valószínűséggel tudja rá a helyes választ, de ha nem tudja, akkor $\dfrac{1}{3}$ valószínűséggel tippel.

Tegyük fel, hogy az illető a kérdésünkre helyes választ adott.
Mennyi a valószínűsége, hogy tudta a választ?

$T$ esemény: tudja a helyes választ.

$J$ esemény: jó/helyes választ adott.

Általánosítsuk a számítást arra az esetre, hogy ha az illető $p$ valószínűséggel tudja a választ!
Vizsgáljuk azt meg, hogy mennyi alternatív válaszlehetőségre lenne szükség akkor, hogy $0.9$ valószínűséggel kijelenthető legyen, hogy az illető azért válaszolt jól, mert tudta a helyes megoldást?

\hrule

\medskip

\noindent \textbf{Megoldás}

\noindent A probléma többféleképpen ábrázolható.

\noindent \textbf{1. változat}

Tegyük fel, hogy a jó választ $J$-vel, a két rossz választ $R_1$-el és $R_2$-vel jelöljük.
Jelölje $T$ azt, hogy ha tudjuk a helyes választ.
Tekintsük a döntést egy mintavételnek, amely esetében a három eshetőség ($J$, $R_1$ és $R_2$) közül választunk.
A megoldás ismeretében csak a helyes választ választjuk, így a helytelen válaszok esetei is átkerülhetnek ahhoz.
Így az alábbi táblázatot kapjuk.

\medskip

\begin{center}
\begin{tabular}{l|c|c|c|}
& $J$ & $R_1$ & $R_2$ \\
\hline
$T$ & 3 & 0 & 0 \\
\hline
$\overline{T}$ & 1 & 1 & 1 \\
\hline
\end{tabular}
\end{center}

\medskip

Összesen 6 választható elemünk van.
\begin{itemize}
\item Annak a valószínűsége, hogy tudjuk a megoldást és helyesen válaszolunk az $P(T \cap J) = \dfrac{3}{6}$.
\item Annak a valószínűsége, hogy jó választ adunk az
\[
P(J) = P(J|T) \cdot P(T) + P(J|\overline{T}) \cdot P(\overline{T})
= 1 \cdot \dfrac{1}{2} + \dfrac{1}{3} \cdot \dfrac{1}{2}
= \dfrac{2}{3}.
\]
\item Annak a valószínűsége, hogy tudjuk a megoldást, feltéve hogy jó választ adtunk:
\[
P(T|J) = \dfrac{P(T \cap J)}{P(J)} =
\dfrac{\frac{1}{2}}{\frac{2}{3}} =
\dfrac{3}{4}.
\]
\end{itemize}

\noindent \textbf{2. változat}

Tegyük fel, hogy az eseteket úgy írjuk fel táblázatos formában, hogy az egyes sorokban az adott eseményhez tartozó feltételes valószínűségek szerepelnek. (Ekkor megkötés, hogy a sorokban lévő értékek összege 1 legyen.)

\medskip

\begin{center}
\begin{tabular}{l|c|c|}
& $J$ & $\overline{J}$ \\
\hline
$T$ & 1 & 0 \\
\hline
$\overline{T}$ & $\dfrac{1}{3}$ & $\dfrac{2}{3}$ \\
\hline
\end{tabular}
\end{center}

\medskip

A számítások az előző változathoz adódnak.

\noindent \textbf{3. változat}

\medskip

Megadhatjuk táblázatos formában közvetlenül az események valószínűségét. Ekkor az alábbit kapjuk.

\medskip

\begin{center}
\begin{tabular}{l|c|c|}
& $J$ & $\overline{J}$ \\
\hline
$T$ & $\dfrac{1}{2}$ & 0 \\
\hline
$\overline{T}$ & $\dfrac{1}{6}$ & $\dfrac{2}{6}$ \\
\hline
\end{tabular}
\end{center}

\medskip

Ekkor szükségszerű, hogy a táblázatban lévő értékek összege 1 legyen (a teljes valószínűség tétele szerint). Mivel a táblázat celláiban szereplő események diszjunktak, ezért az események valószínűségét a sor és az oszlop összegek is megadják.

\begin{itemize}
\item A $P(\overline{T})$ valószínűséget a második sorban lévő értékek összegéből kapjuk, ami így $\dfrac{1}{2}$.
\item $P(J) = \dfrac{1}{2} + \dfrac{1}{6} = \dfrac{2}{3}, \quad P(\overline{J}) = 0 + \dfrac{2}{6}$.
\item Annak a valószínűsége, hogy tudjuk a megoldást, feltéve hogy jó választ adtunk:
\[
P(T|J) = \dfrac{P(T \cap J)}{P(J)} =
\dfrac{\frac{1}{2}}{\frac{2}{3}} =
\dfrac{3}{4}.
\]
\end{itemize}

\hrule

\subsubsection{Ládikók}

Van 3 kincsesládánk, bennük arany és ezüst érmék.
Az elsőben 5 arany van, és 1 ezüst, a másodikban 3 arany és 4 ezüst, a harmadikban pedig 2 ezüst.

A ládikák közül véletlenszerűen választva mennyi a valószínűsége, hogy arany érmét húztunk?

Amennyiben egy arany érmét vettünk ki, mennyi a valószínűsége, hogy az az első, második vagy a harmadik ládikából származik?

\subsubsection{Gyártók és kategóriák}

4 gyártó 3 féle kategóriájú terméket gyárt az alábbi táblázatnak megfelelően.

\begin{tabular}{l|ccc|}
& $B_1$ & $B_2$ & $B_3$ \\
\hline
$A_1$ & 5 & 2 & 1 \\
$A_2$ & 1 & 0 & 3 \\
$A_3$ & 2 & 4 & 0 \\
$A_4$ & 1 & 0 & 1 \\
\hline
\end{tabular}

$A_i$: A termék az $i$-edik gyártótól származik.

$B_j$: A termék a $j$-edik kategóriába tartozik.

Számítsuk ki a valószínűségeket, feltételes valószínűségeket!

\begin{itemize}
\item Mennyi a valószínűsége, hogy a termék az $i$-edik gyártótól származik?
\item Mennyi a valószínűsége, hogy a $j$-edik kategóriába tartozik?
\item Mennyi az $A_i \cap B_j$ valószínűség az egyes helyeken?
\item Mennyi a $P(A_i|B_j)$ valószínűség és a $P(B_j|A_i)$ valószínűség?
\end{itemize}

\subsection{Szúnyogírtás}

Az első írtásnál a szúnyogok 80\%-a pusztult el.
A másodiknál a megmaradtak 40\%-a, a harmadiknál pedig a 20\%-a.
Mennyi a valószínűsége, hogy
\begin{itemize}
\item egy szúnyog túlélte mind a három írtást?
\item kettőt még túl fog élni, feltéve, hogy az elsőt már túlélte? $P(A_3|A_1) = ?$
\end{itemize}

$A_i$: Az $i$-edik írtást túlélte. $A_1 \subset A_2 \subset A_3$

$B_i$: Pontosan az $i$-edik írtásban pusztult el. $B_i \cap B_j = \emptyset, \forall i, j, i \neq j$

\section{Valószínűségi változó, eloszlásfüggvény}

Soroljuk fel az eloszlásfüggvény tulajdonságait!

\[
\xi: \Omega \rightarrow \mathbb{R}
\]

Írjuk fel a szokott rövidített jelöléseket mellőzve, hogy mi az $F_{\xi}(x)$!
\[
F_{\xi}(x) =
P(\xi < x) =
P(\{\xi < x \}) =
P(\{\omega | \xi(\omega) < x\})
\]

\subsection{Változó értékei}

Egy $\xi$ valószínűségi változó a $-2, 8, 12$ értékeket $0.2$, $0.5$ és $0.3$ valószínűségekkel veszi fel.
Írjuk fel a $\xi$-hez tartozó valószínűségeket, eloszlásfüggvényt és ábrázoljuk!
Számítsuk ki a $\xi$ várható értékét!

\subsection{Lehet-e}

Vizsgáljuk meg, hogy a következők lehetnek-e eloszlásfüggvények!

\[
F_1(x) = \begin{cases}
0, & x < \dfrac{1}{2}, \\
\dfrac{x - 1}{x + 1}, & x \geq \dfrac{1}{2}. \\
\end{cases}
\]

\[
F_2(x) = \begin{cases}
0, & x < 1, \\
\dfrac{x - 1}{x + 1}, & x \geq 1. \\
\end{cases}
\]

\[
F_3(x) = \begin{cases}
0, & x < 1, \\
\dfrac{2x - 1}{x + 1}, & x \geq 1.
\end{cases}
\]

\[
F_4(x) = \begin{cases}
0, & x < 0, \\
\dfrac{x^3}{1 + x^2}, & x \geq 0.
\end{cases}
\]

\subsection{Kocka és érmék}

Dobjunk fel egy kockát és 3 érmét.
Az érméknél az eredmény legyen 0 vagy 1.
Rajzoljuk fel az így megadott valószínűségi változó eloszlásfüggvényét!

\medskip

Adott intervallumokra megnézni a valószínűségeket!

\subsection{Négyzet alakú céltábla}

Egy egységnégyzet alakú céltáblára lövünk úgy, hogy biztosan eltaláljuk.
Adjuk meg és rajzoljuk fel a találatok origótól vett távolságainak eloszlásfüggvényét!
\[
F(x) = \begin{cases}
0, & x < 0 \\
\dfrac{x^2 \pi}{4}, & 0 \leq x \leq 1, \\
\dfrac{1}{2} - \dfrac{arccos\left(\dfrac{1}{x}\right)}{\pi} - \sqrt{x^2 - 1}, & 1 < x \leq \sqrt{2}, \\
1, & x > \sqrt{2}.
\end{cases}
\]

\subsection{Sűrűségfüggvényből eloszlásfüggvény}

Tekintsük az alábbi sűrűségfüggvényt!
\[
f(x) =
\begin{cases}
0, & x \leq 2, \\
\dfrac{A}{(1 - x)^2}, & x > 2.
\end{cases}
\]
\begin{itemize}
\item Mennyi az $A$ értéke?
\item Mekkora valószínűséggel esik $\xi$ a $(2, 3)$ intervallumba?
\item Írjuk és rajzoljuk fel a $\xi$-hez tartozó eloszlásfüggvényt!
\end{itemize}

\subsection{Melyik sűrűségfüggvény?}

Vizsgáljuk meg, hogy az alábbiak közül melyik lehet sűrűségfüggvény!

\begin{itemize}
\item[a.]
\[
f(x) =
\begin{cases}
x + \dfrac{1}{2}, & -1 < x < 1, \\
0, & \text{egyébként}.
\end{cases}
\]
\item[b.]
\[
f(x) =
\begin{cases}
x + \dfrac{1}{2}, & -\dfrac{1}{2} < x < \dfrac{3}{2}, \\
0, & \text{egyébként}.
\end{cases}
\]
\item[c.]
\[
f(x) =
\begin{cases}
x + \dfrac{1}{2}, & 0 < x < 1, \\
0, & \text{egyébként}.
\end{cases}
\]
\item[d.]
\[
f(x) =
\begin{cases}
x, & \text{ ha } 0 < x < 1, \\
2 - x, & \text{ ha } 1 \leq x < 2, \\
0, & \text{különben}.
\end{cases}
\]
\item[e.]
\[
f(x) =
\begin{cases}
1 - |1 - x|, & \text{ ha } 0 < x < 2, \\
0, & \text{különben}.
\end{cases}
\]
\end{itemize}

\subsection{Kocka dobássorozat}

Tegyük fel, hogy egy kockával addig dobunk, amíg 6-os értéket nem kapunk.
Az utolsó dobást is beleszámolva mennyi lesz ezen dobássorozatok várható értéke?

\bigskip

Mennyi lesz a várható érték, hogy ha 2 kockával dobunk, és azt várjuk el, hogy a dobásoknál legyen legalább egy hatos?

\bigskip

Konvergens hatványsor:
\(
1 + x + x^2 + x^3 + \cdots =
\displaystyle \sum_{k=0}^{\infty} x^k =
\dfrac{1}{1 - x}
\)

\subsection{Diszkrét várható érték és szórás}

Tegyük fel, hogy a $\xi$ valószínűségi változó a $-6, -3, 1, 2$ értékeket $\dfrac{1}{8}, \dfrac{3}{8}, \dfrac{1}{4}, \dfrac{1}{2}$ valószínűséggel veszi fel.

Mennyi az $E(\xi)$ és a $D(\xi)$ értéke?

\subsection{Folytonos várható érték és szórás}

Számítsuk ki a következő sűrűségfüggvényekkel megadott eloszlások várható értékét és szórását!

Rajzoljuk fel hozzá az eloszlásfüggvényt!

\[
f(x) = \begin{cases}
2x, & 0 < x < 1, \\
0, & \text{egyébként}. \\
\end{cases}
\]

Milyen valószínűséggel esik az érték a $[0.2, 0.5]$ intervallumba?

\[
f(x) = \begin{cases}
|x|, & -1 < x < 1, \\
0, & \text{egyébként}. \\
\end{cases}
\]

Milyen valószínűséggel esik az érték a $[-0.3, 0.7]$ intervallumba?

\subsection{Nem létező várható érték}

Bizonyítsuk be, hogy az alábbi sűrűségfüggvénnyel megadott eloszlásnak nem létezik várható értéke!

\[
f(x) = \begin{cases}
\dfrac{A}{1 + x^2}, & x > 0, \\
0, & \text{egyébként}.
\end{cases}
\]

\[
\displaystyle \int \! \dfrac{1}{1 + x^2} \; \mathrm{d}x = \text{arctg}(x) + c
\]

\subsection{Egyenletes eloszlás}

Tegyük fel, hogy $\xi \in U[5, 20]$.
$E(\xi) = ?, D(\xi) = ?$

\[
E(\xi) = \dfrac{a + b}{2}, \quad D^2(\xi) = \dfrac{(b - a)^2}{12}
\]

Rajzoljuk fel az eloszlás és sűrűségfüggvényt!

\subsection{Binomiális eloszlás}

Tízszer megpróbálunk beletalálni egy céltáblába, amit $0.23$ valószínűséggel találunk el.
$\xi$ jelölje a találatok számát! $E(\xi) = ?, D(\xi) = ?$

\[
E(\xi) = n \cdot p, \quad D^2(\xi) = n \cdot p \cdot q
\]

Mennyi a valószínűsége, hogy a találatok száma legalább 5, de kevesebb mint 9?

\[
P(\xi = k) = \binom{n}{k} \cdot p^k \cdot (1 - p)^{n - k}
\]

\subsection{Kosárlabdázás}

Alíz és Bob kosárlabdáznak.
Alíz $\dfrac{2}{5}$, Bob pedig $\dfrac{1}{4}$ valószínűséggel dob kosárra.
Alíz kezd dobni, és felváltva dobnak, maximum 4-szer az első kosárig.
Jelölje $\xi$ a dobások számát!

\begin{itemize}
\item Határozzuk meg a $\xi$ eloszlását!
\item Számítsuk ki $\xi$ várható értékét és szórását!
\item Számítsuk ki a $P(\xi < 3), P(1 < \xi), P(1 \leq \xi < 4)$ valószínűségeket!
\end{itemize}

\subsection{Busz}

Egy buszon 40 utasnak van hely.
(A vezetőt nem tekintjük utasnak.)
Tegyük fel, hogy $\dfrac{2}{3}$ annak a valószínűsége, hogy egy helyen ül valaki.

\begin{itemize}
\item Mennyi az utasok várható száma?
\item Mi a valószínűsége, hogy minden hely foglalt?
\item Mennyi a valószínűsége, hogy legalább 4-en utaznak a buszon?
\end{itemize}

\subsection{Selejtes gyártás}

Egy gyártás során a termékek 15\%-a selejtes lesz.
12 elemű mintát veszünk belőle.
Mennyi a valószínűsége annak, hogy
\begin{itemize}
\item a mintában legfeljebb 2 selejtes termék van,
\item 7 jó termék van,
\item legalább 2, de legfeljebb 5 selejtes termék van?
\end{itemize}
Mennyi a mintában lévő selejtek számának várható értéke?

\subsection{Szakaszon pontok}

Adott egy $[0, r]$ szakasz.
Véletlenszerűen kiválasztunk rajta 2 értéket.
Mennyi a valószínűsége, hogy ezen értékek négyzetösszege nagyobb lesz, mint $r^2$?

\subsection{Tanuló és jegyzetelés}

Egy hallgató a tanulásra fordított ideje negyedét matematika, 40\%-át angol, a többi időt pedig egyéb tárgyak tanulásával tölti.
Matematika tanulás közben 0.6, angol tanulás közben 0.4 egyéb tárgyak tanulása közben pedig 0.3 annak a valószínűsége, hogy éppen ír.

\medskip

Azt látjuk, hogy tanulás közben a hallgató éppen ír.
Mennyi a valószínűsége, hogy éppen matematikát tanul?

\subsection{Intervallum meghatározás}

Adott egy eloszlás a következő sűrűségfüggvénnyel!

\[
f(x) = \begin{cases}
\dfrac{1}{x}, & e^5 < x < m, \\
0, & \text{egyébként}.
\end{cases}
\]

\begin{itemize}
\item Határozzuk meg a hiányzó $m$ értéket!
\item Rajzoljuk fel a sűrűség- és az eloszlásfüggvényét!
\item Számítsuk ki a várható értékét és a szórását!
\end{itemize}

\section{Speciális eloszlások}

Poisson eloszlás

\[
P(\xi = k) = \dfrac{\lambda^k}{k!} \cdot e^{-\lambda}
\]

\[
E(\xi) = \lambda, \quad D^2(\xi) = \lambda 
\]

Exponenciális eloszlás

\[
F(x) = \begin{cases}
1 - e^{-\lambda x}, & x \geq 0, \\
0, & x < 0.
\end{cases}, \quad
f(x) = \begin{cases}
\lambda \cdot e^{-\lambda x}, & x \geq 0, \\
0, & x < 0.
\end{cases}
\]

Normális eloszlás

\[
\mathcal{N}(\mu, \sigma^2), \quad
\varphi(x) = \dfrac{1}{\sigma \sqrt{2\pi}} \cdot e^{-\dfrac{(x - \mu)^2}{2\sigma^2}}
\]

\[
\Phi(-x) = 1 - \Phi(x)
\]

Geometriai eloszlás

\[
P(\xi = k) = q^{k - 1} \cdot p, \quad
E(\xi) = \dfrac{1}{p}, \quad
D^2(\xi) = \dfrac{q}{p^2}
\]

\subsection{Csillaghullás}

Egy augusztusi éjszakán átlagosan 10 percenként látunk egy hullócsillagot.
\begin{itemize}
\item Mennyi a valószínűsége, hogy negyedórán belül kettőt is látni fogunk?
\item Láttunk egy hullócsillagot.
Mennyi a valószínűsége, hogy 5 percen belül láthatunk még egyet?
\end{itemize}

\subsection{Kalács}

Kalácssütéskor 1 kg tésztába 30 szem mazsolát tesznek.
Mennyi a valószínűsége, hogy egy 5 dekagramos szeletben kettőnél több mazsolaszem lesz?

% Poisson eloszlás, lambda = 30/20, p = 0.1912

\subsection{Szövet}

Egy szöveten 100 méterenként átlagosan 5 hiba van.
A hibák Poisson eloszlással fordulhatnak elő.
Egy 300 méteres szövetet 3 méteres darabokra vágnak.

\medskip

Előreláthatólag mennyi hibátlan darab lesz közöttük?

\subsection{Rendelés}

Kínából alkatrészeket rendelünk.
A tapasztalatok alapján az alkatrészek általában 7 nap alatt érkeznek meg.
A kiszállítás idejének eloszlása exponenciális.

\medskip

\begin{itemize}
\item Mennyi a valószínűsége, hogy egy alkatrész 3 napon belül megérkezik?
\item Mennyi a valószínűsége, hogy egy alkatrészre 5 napnál többet kell várni?
\end{itemize}

\subsection{Programhiba}

Egy programot elindítva az első hiba előfordulásáig $\xi$ idő telik el, ami a megfigyelések alapján exponenciális eloszlású és 6 óra várható értékű.

\medskip

Mennyi a valószínűsége, hogy a programban 8 órán belül nem jelentkezik hiba?

\subsection{Benzinkút}

Egy benzinkútnál annak a valószínűsége, hogy a tankolásnál 6 percnél többet kell várni a tapasztalatok alapján 0.1.
A várakozási idő exponenciális eloszlású.

\medskip

Mennyi a valószínűsége, hogy véletlenszerű időpontban a benzinkúthoz érve 3 percen belül sorra kerülünk?

\subsection{Bolt}

Egy forgalmas boltba a vevők Poisson eloszlás szerint érkeznek, átlagosan 40 vevő óránként.

\medskip

Mennyi a valószínűsége, hogy egy vevő érkezése után 5 percig nem érkezik másik vevő?

\subsection{Állatos}

Egy kisállat várható élettartama exponenciális eloszlás szerint változik 12 év várható élettartammal.

\medskip

Mennyi a valószínűsége, hogy a kisállat az életének 8. évében fog elpusztulni?


\subsection{Intervallum becslés}

Adott egy $\xi$ valószínűségi változó, amelynek 3 a várható értéke és 2 a szórása.
Mekkora legyen az $A$ érték, hogy ha azt szeretnénk, hogy a $\xi$ értéke legalább $\dfrac{1}{2}$-ed valószínűséggel essen a $(2, A)$ intervallumba?

% Legalább 4.74-nek kell lennie.

\subsection{Munkapad}

A munkapadról kikerült termékek hossza normális eloszlású valószínűségi változó, 20cm várható értékkel és 0.2cm szórással.
\begin{itemize}
\item Mekkora a valószínűsége, hogy egy termék hossza 19,7 és 20.3 cm közé fog esni?
\item Milyen pontosságot biztosíthatunk 0.95 valószínűséggel a munkadarabok hosszára?
\end{itemize}

\subsection{Szolgáltató, megrendelése}

Egy szolgáltatóhoz a naponta érkező megrendelések száma közelítőleg normális eloszlású, és tudjuk, hogy 
szórása $\sigma = 10$.
Mekkora a megrendelések várható száma, hogy ha tudjuk, hogy $P(\xi < 20) = 0.1$.

\subsection{Binomiális eloszlás közelítése}

Elvégzünk 1000 kísérletet, és tudjuk, hogy egy kisérlet 0.27 valószínűséggel lehet sikeres.
Jelölje $\xi$ a sikeres kísérletek számát.

\medskip

Adjunk közelítést annak a valószínűségére vonatkozóan, hogy a $\xi$ a $[250, 280]$ intervallumba fog esni.

\subsection{Diótörés}

Diót törünk, és azt tapasztaljuk, hogy a diók 0.88 valószínűséggel lesznek jók.

\medskip

\begin{itemize}
\item Fogalmazzuk meg erre vonatkozóan, hogy mit vizsgál a binomiális eloszlás és a geometriai eloszlás.
\item Adjunk rá konkrét példát várható értékkel és szórással!
\end{itemize}


\section{Egyenlőtlenségek}

\subsection{Markov}

Tudjuk, hogy egy $\xi$ valószínűségi változó sűrűségfüggvénye csak a $(0, 6)$ intervallumba nem 0, illetve, hogy $E(\xi) = 1$.
Lássuk be, hogy a $P(\xi < 5) \geq \dfrac{4}{5}$ teljesül!

\medskip

\textbf{Markov egyenlőtlenség}

\[
P(\xi \geq c) \leq \dfrac{E(\xi)}{c}
\]

\subsection{Poission, Csebisev}

Egy újságárus óránként eladott újságainak száma Poisson eloszlású $\lambda = 64$ paraméterrel.
Adjunk alsó becslést a $P(48 < \xi < 80)$ értékére a Csebisev-egyenlőtlenség segítségével!

\medskip

\textbf{Csebisev egyenlőtlenség}

\[
P(|\xi - E(\xi)| \geq \varepsilon) \leq \dfrac{D^2(\xi)}{\varepsilon^2}
\]

\subsection{Nagy számok törvénye}

Egy célpontra 200 lövést adunk le.
A találatok valószínűsége minden lövésnél 0.4.
Milyen határok közé fog esni 0.9 valószínűséggel a találatok száma?

\medskip

\textbf{Bernoulli-féle nagy számok törvénye}

\[
P\left(\left|\dfrac{\xi_{n}}{n} - p\right| \geq \varepsilon\right)
\leq \dfrac{p \cdot q}{n \cdot \varepsilon^2}
\leq \delta
\]

\subsection{Ismeretlen valószínűség}

Egy csavargyárban 5000 csavart megvizsgálva 80 selejteset találtunk.
90\%-os valószínűséggel milyen intervallumon belül lesz annak a valószínűsége, hogy egy elkészült csavar selejt?

\medskip

\begin{itemize}
\item Egy Bernoulli-kísérletsorozatról van szó.
\item Írjuk fel a nagy számok törvényét!
\item Adjunk felső becslést az egyenlőség jobb oldalára vonatkozóan!
\item Ábrázoljuk hozzá a $p(1 - p)$ függvényt a $[0, 1]$ intervallumon!
\end{itemize}

\[
P\left(\left|\dfrac{\xi_{n}}{n} - p\right| \geq \varepsilon\right)
\leq \dfrac{p \cdot q}{n \cdot \varepsilon^2}
\leq \dfrac{1}{4 n \varepsilon^2}
\leq \delta
\]

\section{Véletlen vektorok}

\subsection{Eloszlás vizsgálata}

Adottak $(\xi, \eta)$ értékei az alábbi táblázatnak megfelelően.

\begin{tabular}{l|c|c|}
$(\xi, \eta)$ & 0 & 1 \\
\hline
0 & $p$ & $p$ \\
1 & $p$ & $3p$ \\
2 & $2p$ & $4p$ \\
\hline
\end{tabular}

$p = \dfrac{1}{12}$

\begin{itemize}
\item Írjuk és rajzoljuk fel az eloszlásfüggvényt!
\item Számítsuk ki a $P(\xi = i|\eta = 0), P(\xi < 2|\eta = 0), P(\xi \geq 1|\eta = 1), P(\eta = 1|\xi \geq 1)$ értékeket!
\end{itemize}

\begin{solution}
\[
P(\xi = 0|\eta = 0) = \dfrac{P(\xi = 0, \eta = 0)}{P(\eta = 0)} = \ldots
\]
\end{solution}

\subsection{Egyenletes}

Adott egy $(\xi, \eta)$ valószínűségi vektorváltozó a $(-a \leq x \leq a, -b \leq y \leq b)$ egyenlőtlenségekkel jellemzett téglalap alakú tartomány felett, amely a tartomány felett egyenletes eloszlású.

\begin{itemize}
\item Írjuk fel $(\xi, \eta)$ eloszlásfüggvényét!
\item Írjuk fel a peremeloszlás függvényeket is!
\item Számítsuk ki az alábbi valószínűségeket az eloszlásfüggvény alapján!
\end{itemize}
\[
P(\xi < 0, \eta < 0), P\left(\xi < 0, \eta \geq \dfrac{b}{2}\right),
P\left(0 \leq \xi < \dfrac{a}{2}, 0 \leq \eta \leq \dfrac{b}{2}\right), P\left(\xi > \dfrac{a}{2}\right)
\]

\begin{solution}
Felrajzolni az eloszlásfüggvényt "3D-ben"!

\[
P(x_1 \leq \xi \leq x_2, y_1 \leq \eta \leq y_2) =
F(x_2, y_2) - F(x_1, y_2) - F(x_2, y_1) + F(x_1, y_2)
\]
\end{solution}

\subsection{Egyenletes 2}

Legyen $(\xi, \eta)$ egyenletes eloszlású a $0 \leq x \leq a, 0 \leq y \leq a - x$ tartományon.

\begin{itemize}
\item Írjuk fel az eloszlásfüggvényt!
\item Számítsuk ki a $\xi$ és az $\eta$ peremeloszlásfüggvényét!
\end{itemize}

\begin{solution}
\begin{itemize}
\item Fel kell rajzolni a tartományt!
\item Meg kell határozni $a$ értékét!
\item Fel kell írni az eloszlásfüggvényét!
\item Fel kell írni a peremeloszlásfüggvényeket!
\item Fel kell írni az együttes sűrűségfüggvényt!
\item Integrálással (az együttes sűrűségfüggvényből) meg kell határozni a peremsűrűségfüggvényeket!
\end{itemize}
\end{solution}

\subsection{Általános}

\[
f(x, y) = \begin{cases}
e^{-x-y}, & x > 0, y > 0, \\
0, & \text{egyébként}.
\end{cases}
\]

\[
P(\xi < 1, \eta < 1), P\left(\xi < 1, \eta \geq \dfrac{3}{2}\right)
\]

% Nézzünk meg háromszög feletti integrált is!

\begin{solution}
\begin{itemize}
\item Ellenőrízzük, hogy valóban sűrűségfüggvény!
\item Figyelni kell, hogy a függvény csak a pozitív síknegyedben van értelmezve!
\end{itemize}
\end{solution}

\subsection{Négyzetes}

\[
f(x, y) = x^2 + A y^2, \quad (0 < x < 1, 0 < y < 2), \quad A = ?
\]

Írjuk fel a peremeloszlás függvényeket!

\begin{solution}
\begin{itemize}
\item Figyelni kell a $(-\infty, +\infty)$ intervallum kezelésére!
\end{itemize}
\end{solution}

\subsection{Feltételes eloszlásfüggvény}

\[
f(x, y) = \begin{cases}
A, & x \geq 0, y \geq 0, y \leq \dfrac{3 - x}{2}, y \leq -2x + 3, \\
0, & \text{különben}.
\end{cases}
\]

\begin{itemize}
\item Határozzuk meg az $A$ értékét!
\item Írjuk fel a peremeloszlásfüggvényeket!
\end{itemize}

\[
P\left(\xi < x, 1 < \eta < \dfrac{3}{2}\right) = ?
\]

\[
F\left(x|1 < \eta < \dfrac{3}{2}\right) = ?
\]

\[
F(x|1 < \eta < 3) = ?
\]

\begin{solution}
\begin{itemize}
\item Felrajzolni jó nagy ábrára, hogy látszódjanak a háromszögek!
\item Felírni az eseteket!
\end{itemize}
\end{solution}

\subsection{Függetlenség}

\[
f(x, y) = \begin{cases}
1, & 0 < x < 1, 0 < y < 2(1 - x), \\
0, & \text{egyébként}. \\
\end{cases}
\]

\begin{itemize}
\item Vizsgáljuk meg, hogy $\xi$ és $\eta$ függetlenek-e!
\item Írjuk fel az $f(x|y)$ és $f(y|x)$ feltételes sűrűségfüggvényeket!
\item Számítsuk ki a $P\left(0 < \xi < \dfrac{1}{2}\mid\eta = \dfrac{1}{2}\right)$ értékét!
\end{itemize}

\begin{solution}
\[
f_{\xi}(x) = ?, f_{\eta}(y) = ?, f(x|y) = \dfrac{f(x, y)}{f(y)}, \ldots
\]
\end{solution}

\subsection{Összeg}

\[
f_{\xi\eta} =
\begin{cases}
e^{-x-y}, & x \geq 0, y \geq 0, \\
0, & \text{egyébként}.
\end{cases}
\]

\[
f_{\xi + \eta}(z) = f_{\zeta}(z) = ?
\]

\subsection{Összeg és hányados}

\[
f(x, y) = \begin{cases}
\dfrac{1}{4}, & 0 < x < 2, 0 < y < 2, \\
0, & \text{egyébként}. \\
\end{cases}
\]

\[
\zeta = \xi + \eta, \tau = \dfrac{\xi}{\eta}
\]

\[
f_{\xi + \eta}(z) = ?, \quad f_{\tau}(z) = ?
\]

\subsection{Várható érték}

\[
f(x, y) = \begin{cases}
x + y, & 0 < x < 1, 0 < y < 1, \\
0, & \text{egyébként}. \\
\end{cases}
\]

\[
E(\xi + \eta) = ?, E(\xi \cdot \eta) = ?
\]

\subsection{Korrelációs együttható}

\begin{tabular}{|l|cc|}
\hline
$(\xi, \eta)$ & 1 & 0 \\
\hline
1 & 0.5 & 0.04 \\
0 & 0.06 & 0.4 \\
\hline
\end{tabular}

Számítsuk ki $\xi$ és $\eta$ korrelációs együtthatóját!

\subsection{Regressziós függvények}

\begin{tabular}{l|c|c|c|}
$(\xi, \eta)$ & 0 & 1 & 3 \\
\hline
0 & $\dfrac{3}{20}$ & $\dfrac{6}{20}$ & $\dfrac{4}{20}$ \\
2 & $\dfrac{1}{20}$ & $\dfrac{5}{20}$ & $\dfrac{1}{20}$ \\
\hline
\end{tabular}

Határozzuk meg a $\xi$-nek az $\eta$-ra vonatkozó, és az $\eta$-nak a $\xi$-re vonatozó regressziós függvényét!

\begin{solution}
\[
q_i = P(\xi = x_i), r_j = P(\eta = y_j)
\]

\[
E(\xi|\eta = y_j) = \displaystyle\sum_i x_i \dfrac{p_{ij}}{r_j} = m_2(y_j)
\]

\[
E(\eta|\xi = x_i) = \displaystyle\sum_j y_j \dfrac{p_{ij}}{q_i} = m_1(x_i)
\]

Elég csak jelölni, hogy mik a $p_{ij}, q_i, r_j$ értékek, és szépen sorban visszahelyettesíteni!
\end{solution}

\end{document}
